\section{总结与展望}\label{chap:conclusion}

\subsection{工作总结}

本文围绕可证安全图像隐写系统的设计与实现展开研究，主要完成了以下工作：

\begin{enumerate}
    \item \textbf{双算法隐写引擎}：集成了 Pulsar 和 SparSamp 两种基于扩散模型的可证安全隐写算法，设计了统一的算法注册与切换机制。Pulsar 基于像素空间扩散模型（DDIM 采样器），通过控制最后一步调度器的方差噪声来源和变码率级联纠错码实现消息嵌入；SparSamp 基于稀疏采样思想，在 IDDPM P2-weighting 框架下通过消息驱动采样和稀疏采样消歧机制实现消息嵌入。两种算法共享相同的 API 接口，支持运行时切换。SparSamp 实现了状态缓存机制，避免重复的扩散采样计算。

    \item \textbf{端到端加密机制}：实现了基于对称助记词的端到端加密方案。用户通过助记词短语经 PBKDF2-SHA256（100,000 次迭代）派生 256 位用户密钥；通信双方的用户密钥按字典序排列后经 HKDF-SHA256 协商出对称消息密钥；消息使用 AES-256-GCM 加密封装为 SealedMessage 格式。消息密钥仅在内存中存在，不进行持久化存储。在隐写聊天模式下，加密后的密文作为隐写嵌入的载荷，实现了"先加密、再隐写"的双重保护。

    \item \textbf{系统架构设计}：设计了基于客户端--服务端的跨平台隐写通信系统架构。服务端基于 Python FastAPI 框架构建，集成双算法隐写引擎和 WebSocket 实时通信服务；客户端分别使用 Kotlin Jetpack Compose（Android）和 Vue.js 3（Web）开发，实现了完整的用户交互、本地数据管理和端到端加解密。

    \item \textbf{实时通信系统}：基于 WebSocket 协议实现了全双工实时通信，支持消息即时推送、完整的消息生命周期管理（已发送/已送达/已读回执/消息撤回）、在线状态感知、好友请求通知和离线消息缓存。设计并实现了单端登录安全策略，防止多设备同时登录导致的数据一致性问题。

    \item \textbf{用户体系与社交功能}：实现了完整的用户注册登录系统（PBKDF2 密码哈希 + JWT 认证）、基于邀请码的好友添加机制，以及好友请求的异步处理流程。支持联系人屏蔽功能，被屏蔽用户的后续消息自动过滤。

    \item \textbf{隐写通信功能}：在聊天界面中集成了隐写模式，用户可在普通消息和隐写消息之间无缝切换。系统自动使用通信双方的邀请码 XOR 运算派生隐写密钥，实现了端到端的隐蔽信息传输。同时提供了独立的隐写工具页面，支持手动嵌入和提取操作，支持算法和模型选择。

    \item \textbf{多平台客户端开发}：Android 客户端采用 Material 3 Expressive 薄荷色主题，包含 9 个页面和 4 个 ViewModel；Web 客户端采用 Aegis Slate 暗色 Material 3 设计系统，包含 10 个页面组件和 4 个 Pinia Store。两个客户端的功能完全对齐，支持跨平台互通。

    \item \textbf{系统测试验证}：对系统进行了全面的功能测试、隐写性能分析和跨平台兼容性测试，共计 60 个测试用例全部通过，覆盖了用户管理、好友管理、消息生命周期、端到端加密、双算法隐写和跨平台互操作等核心功能。
\end{enumerate}

\subsection{不足与展望}\label{sec:future_work}

尽管本系统已实现了基本完整的可证安全隐写通信功能，但仍存在一些不足之处，有待在后续工作中改进：

\begin{enumerate}
    \item \textbf{嵌入容量有限}：受限于扩散模型的采样步数和图像分辨率，当前系统的单幅图像嵌入容量约为数百字节，仅适合短文本消息的传输。未来可通过支持更大分辨率的模型、优化编码算法或采用多图像分块嵌入策略来提升容量。

    \item \textbf{计算延迟较高}：两种隐写算法均涉及多步扩散模型推理，单次嵌入或提取操作耗时较长。SparSamp 通过状态缓存机制缓解了重复计算的问题，但首次请求仍需约 50 秒。未来可通过模型量化、蒸馏或使用更高效的扩散模型采样器来降低延迟。

    \item \textbf{端到端加密方案的局限性}：当前采用的对称助记词方案缺乏前向安全性和未来安全性。未来可考虑引入 Signal 协议的双棘轮算法，实现密钥的自动轮换和前向安全保护。

    \item \textbf{隐写密钥派生方案的安全性不足}：当前隐写密钥通过双方邀请码 XOR 运算生成（$K_{\text{stego}} = \text{inviteCode}_A \oplus \text{inviteCode}_B$），该方案存在明显局限：邀请码仅为 8 位字母数字字符（约 48 比特熵），XOR 运算不增加密钥空间；且邀请码存储在服务端数据库中，服务端管理员或数据库泄露均可推导出隐写密钥。这意味着隐写密钥的安全性依赖于服务端的可信度，与"可证安全"的理想目标存在差距。未来可考虑基于双方助记词通过 HKDF 派生独立的隐写密钥，使其安全性与端到端加密密钥保持一致。

    \item \textbf{服务端集中式计算}：当前隐写计算完全在服务端执行，存在单点瓶颈和隐私风险。未来可探索将部分计算迁移到客户端（如使用 WebAssembly 或移动端 GPU 加速），或采用分布式计算架构分担负载。

    \item \textbf{群组聊天支持}：当前系统仅支持一对一聊天。未来可扩展为群组聊天功能，但需要解决群组密钥管理和多方隐写密钥协商的问题。

    \item \textbf{好友关系本地存储的局限性}：当前系统的好友关系仅存储在客户端本地数据库中，服务端不维护好友列表。这意味着用户更换设备或清除本地数据后，好友关系将丢失，需要重新添加。未来可在服务端引入加密存储的好友关系表，或利用端到端加密的助记词在新设备上恢复好友关系。

    \item \textbf{文件类型扩展}：当前系统仅支持文本消息的隐写传输。未来可扩展支持文件、图片等多种数据类型的隐写嵌入，拓展应用场景。
\end{enumerate}

总体而言，本文设计并实现的可证安全图像隐写系统将可证安全隐写理论与端到端加密技术相结合，通过双算法支持和跨平台实现，验证了可证安全隐写技术在实际通信系统中的可行性。系统在功能完整性、跨平台兼容性和通信安全性方面均达到了预期目标，60 个测试用例的通过证明了所设计架构的正确性和实用性。尽管在嵌入容量和计算效率方面仍有提升空间，但本文的工作表明，基于扩散模型的可证安全隐写技术已具备在实际通信场景中部署的基本条件，为安全隐蔽通信领域从理论走向工程应用提供了一种可参考的实现方案。
